\begin{theorem}[7.7.2]
\EGAReaderTarget{ega4.c04.thm.7-7-2}{IV.2.s7.7.2-fr}
(Nagata). --- Sea $A$ un anillo noetheriano. Las condiciones siguientes son
equivalentes:
\begin{enumerate}
\item[\emph{a)}] $A$ es un anillo universalmente japonés.
\item[\emph{b)}] Todo anillo cociente de $A$ que sea un dominio de integridad
es un anillo japonés.
\item[\emph{c)}] Para todo ideal maximal $\mathfrak m$ de $A$, todo anillo
cociente de $A_{\mathfrak m}$ que sea un dominio de integridad es un anillo
japonés y, para todo anillo cociente $B$ de $A$ que sea un dominio de
integridad, con cuerpo de fracciones $K$, toda extensión finita puramente
inseparable $K'$ de $K$ y toda sub-$B$-álgebra finita $B'$ de $K'$, con
cuerpo de fracciones $K'$, existe $f'\neq0$ en $B'$ tal que $B'_{f'}$ es
integralmente cerrado \textup{(cf. (6.13.7))}.
\end{enumerate}

Además, si $A$ satisface estas condiciones, también las satisfacen toda
localización $S^{-1}A$ de $A$ y toda $A$-álgebra de tipo finito.
\end{theorem}

\begin{proof}
Demostremos primero que $b)$ implica $c)$; todo anillo cociente de
$A_{\mathfrak m}$ que sea un dominio de integridad es una localización de
un anillo cociente de $A$ que sea un dominio de integridad y, por tanto,
un anillo japonés (\textbf 0, 23.1.1). Por otra parte, todo cociente $B$ de $A$ por
uno de sus ideales primos es un anillo japonés, y lo mismo ocurre con $B'$
(\textbf 0, 23.1.1). De aquí se deduce que el conjunto
$\operatorname{Nor}(\operatorname{Spec}(B'))$ es abierto (6.13.3) y, por
tanto, contiene un abierto no vacío
$D(f')=\operatorname{Spec}(B'_{f'})$.

\oldpage[IV-2]{213}

Demostremos, en segundo lugar, que, si $A$ satisface $c)$, también la
satisface toda $A$-álgebra de tipo finito $B$. En primer lugar, para todo
ideal primo $\mathfrak q$ de $B$, $B_{\mathfrak q}$ es el anillo local en
un ideal primo de $S^{-1}B$, donde $S=A\setmin\mathfrak p$, siendo
$\mathfrak p$ la imagen inversa de $\mathfrak q$ en $A$; como, por
hipótesis, todo anillo cociente de $A_{\mathfrak p}$ que sea un dominio
de integridad es un anillo japonés, lo mismo ocurre con todo anillo
cociente de $B_{\mathfrak q}$ que sea un dominio de integridad, pues
$S^{-1}B$ es una $A_{\mathfrak p}$-álgebra de tipo finito (7.6.5).
Por otra parte, si $C$ es un cociente de $B$ que es un dominio de
integridad, con cuerpo de fracciones $K$, $K'$ es una extensión finita
puramente inseparable de $K$ y $C'$ es una sub-$C$-álgebra finita de
$K'$, con cuerpo de fracciones $K'$, entonces $C'$ es una $A$-álgebra de
tipo finito que es un dominio de integridad. En virtud de (6.13.7), la
hipótesis $c)$ sobre $A$ implica que el espectro de $C'$ contiene un
abierto no vacío cuyos puntos son todos normales; dicho de otro modo,
existe $g'\neq0$ en $C'$ tal que $C'_{g'}$ es integralmente cerrado.
Esto demuestra nuestra afirmación.

Este resultado muestra que, para demostrar la equivalencia de $a)$, $b)$
y $c)$, basta demostrar que $c)$ implica $b)$, e incluso basta probar que,
para todo anillo cociente $B$ de $A$ que sea un dominio de integridad,
con cuerpo de fracciones $K$, el cierre entero $B'$ es un $B$-módulo
finitamente generado. Ahora bien, la condición $c)$ implica que existe
$f\neq0$ en $B$ tal que $B_f$ es integralmente cerrado; de este modo se
ha verificado para $B$ la condición (i) de (6.13.6). Por otra parte, $B$
satisface también la condición (ii) de (6.13.6). En efecto, para todo
ideal maximal $\mathfrak q$ de $B$, $B_{\mathfrak q}$ es un cociente de
un anillo local $A_{\mathfrak p}$, donde $\mathfrak p$ es maximal en
$A$; como, por hipótesis, todo anillo cociente de $A_{\mathfrak p}$ que sea un dominio de integridad es un anillo japonés,
lo mismo ocurre con $B_{\mathfrak q}$. Además, para todo ideal primo
$\mathfrak r$ de $B$, $B_{\mathfrak r}$ es una localización de un
$B_{\mathfrak q}$, donde $\mathfrak q$ es un ideal maximal de $B$, y es,
por tanto, también un anillo japonés. Esto demuestra nuestra afirmación.

Se ha visto además que, si $A$ satisface las condiciones equivalentes
$a)$, $b)$ y $c)$, también las satisface toda $A$-álgebra de tipo finito.
Por otra parte, si $A$ satisface $b)$, también la satisface toda
localización $S^{-1}A$, pues todo cociente de $S^{-1}A$ que sea un dominio
de integridad tiene la forma $S^{-1}A/S^{-1}\mathfrak p$, donde
$\mathfrak p$ es un ideal primo de $A$; como, por hipótesis,
$A/\mathfrak p$ es un anillo japonés, lo mismo ocurre con
$S^{-1}(A/\mathfrak p)$ (\textbf 0, 23.1.1). Con ello queda demostrada la última
afirmación del enunciado.
\end{proof}
